% Subsection: what the network-contention replay measures, and why CGS loses ground under it.
% Every number is a macro generated by paper/why_contention.py (outputs/tables/contention_macros.tex) and
% checked by paper/verify.py; the table body is outputs/tables/tab_contention_rows.tex.
% Copy outputs/tables/ next to the manuscript as tables/, then \input this file.
\input{tables/contention_macros.tex}

\subsection{Behaviour under network contention}
\label{sec:contention}

All compared methods, CGS included, assume that any number of messages can be in transit at once. To see how
their schedules fare when this does not hold, we replay each schedule under a one-port, full-duplex model:
a processor sends at most one message and receives at most one message at a time, and a message between two
processors occupies the sender's out-port and the receiver's in-port for its whole cost. The replay keeps each
schedule's processor mapping and the order of tasks on every processor, so it measures how a fixed schedule
degrades, not how a contention-aware scheduler would perform.

Under this replay CGS loses its lead on the workflows: its mean paired gain over AEFT is
\WCWfAeftGainCont\,\%, against \WCWfAeftGainPlan\,\% without contention. The difference reflects what the replay measures rather than the
planned length of the schedules.

\paragraph{The replay is dominated by the busiest port.}
Every method's schedules lengthen several-fold under contention: on average by a factor between
\WCFftInflMin{} and \WCFftInflMax{} on FFT and between \WCMontageInflMin{} and \WCMontageInflMax{} on Montage,
depending on the method. Since a port carries one message at a time, the total cost of the messages through the
busiest port is a lower bound on the contention makespan. On the workflows this port load alone exceeds the
planned makespan in \WCWfCgsPortExceeds\,\% of CGS's schedules and \WCWfAeftPortExceeds\,\% of AEFT's, and the
contention makespan is a median \WCWfCgsContOverPort{} times the port load for CGS and \WCWfAeftContOverPort{}
times for AEFT. Across the nine schedules of the same instance, the logarithm of
the contention makespan has a median correlation of \WCWfCorrPort{} with that of the port load and of
\WCWfCorrPlan{} with that of the planned makespan. The replay therefore ranks schedules by how much
communication their busiest port carries, almost independently of their planned length.

\paragraph{CGS trades a little communication for shorter schedules.}
Relative to AEFT, CGS's workflow schedules move \WCWfCgsVol\,\% more data between processors, which makes their
busiest port \WCWfCgsPort\,\% heavier and their contention makespan \WCWfCgsCont\,\% longer
(Table~\ref{tab:contention}). The busiest port decides the outcome instance by instance: CGS is longer than AEFT
under contention on \WCWfLose{} of the \WCWfN{} test workflow instances, and on \WCWfLoseHeavy\,\% of these its
busiest port is the heavier one; when its port is heavier it is longer in \WCWfLoseIfHeavy\,\% of instances,
when lighter in \WCWfLoseIfLight\,\%. The same trade-off orders the baselines: PEFT and BPPTS plan schedules
\WCWfPeftPlanned\,\% and \WCWfBpptsPlanned\,\% longer than AEFT's, but their communication volume differs from
AEFT's by \WCWfPeftVol\,\% and \WCWfBpptsVol\,\%, and their contention makespan by \WCWfPeftCont\,\% and
\WCWfBpptsCont\,\%. Across all methods in Table~\ref{tab:contention}, the difference in contention makespan follows the
difference in busiest-port load far more closely than the difference in planned makespan.

\paragraph{Joint placement is not the cause.}
CGS without joint placement already moves \WCWfCgsSingleVol\,\% more data than AEFT, so the extra volume comes
from processor selection with the downstream table, not from the closures; forming closures at every
opportunity (no gate) instead changes it to \WCWfCgsNogateVol\,\% relative to AEFT. Nor does CGS concentrate messages on
particular processors: the median ratio of the busiest in-port's load to the mean is \WCWfCgsInConc{} for CGS
and \WCWfAeftInConc{} for AEFT. On the random DAGs at CCR~10, where CGS's communication volume differs from
AEFT's by \WCRHighCgsVol\,\%, its gain over AEFT grows under contention, from \WCRHighAeftGainPlan\,\% to \WCRHighAeftGainCont\,\%.

\paragraph{Implications.}
None of the compared methods models contention, and on these workflows all of them degrade by large factors
when it is present; contention-aware scheduling is a separate problem that we do not address. Within the
contention-free model that CGS and all baselines share, CGS's shorter schedules come with slightly more
communication, and under a one-port network that communication, not the schedule's planned length, sets the
makespan.

\begin{table*}[t]
\centering
\caption{Contention replay on the test split. Inflation: mean ratio of contention to planned makespan on the
workflows. Other columns: mean difference from AEFT (\%, positive = larger than AEFT) in planned makespan,
volume of communication between processors, load of the busiest port (a lower bound on the contention
makespan) and contention makespan, on the workflows and on the random DAGs with CCR~10.}
\label{tab:contention}
\footnotesize
\setlength{\tabcolsep}{4.5pt}
\begin{tabular}{lrrrrrrr}
\toprule
 & & \multicolumn{4}{c}{Workflows: difference from AEFT (\%)} & \multicolumn{2}{c}{Random, CCR 10 (\%)} \\
\cmidrule(lr){3-6}\cmidrule(lr){7-8}
Method & Inflation ($\times$) & Planned & Volume & Busiest port & Contention & Volume & Contention \\
\midrule
\csname @@input\endcsname tables/tab_contention_rows.tex % primitive input: \input is not expandable inside tabular
\bottomrule
\end{tabular}
\end{table*}
